\section{Formal Conjectures Details}\label{sec:app_formal_conjectures}


\subsection{Source and AMS Classification Tables}\label{sec:app_source_ams_tables}

Note that a single source problem typically gives rise to multiple formal statements due to variants (e.g., special cases, generalizations, solved sub-problems, and test lemmas). We explicitly allow and encourage the inclusion of multiple alternative formalizations for a single informal claim to capture different mathematical perspectives or resolutions of ambiguity; thus, the counts in this table reflect formal statements, not distinct source problems.

\begin{table}[h]
    \caption{Number of formal statements by source collection, split by category. A single source problem may yield multiple statements (variants, special cases, tests).}
    \label{tab:problems-by-collection}
    \centering
    \begin{tabular}{lrrrr}
    \toprule
    \textbf{Source Collection} & \textbf{Total} & \textbf{Research Open} & \textbf{Research Solved} & \textbf{All Other} \\
    \midrule
    Erd\H{o}s Problems & 1318 & 551 & 551 & 216 \\
    Wikipedia & 476 & 208 & 113 & 155 \\
    Green's Open Problems & 175 & 83 & 73 & 19 \\
    Papers & 131 & 65 & 25 & 41 \\
    Open Quantum Problems & 125 & 35 & 13 & 77 \\
    Written on the Wall II & 110 & 13 & 12 & 85 \\
    OEIS & 105 & 22 & 6 & 77 \\
    arXiv & 56 & 13 & 13 & 30 \\
    MathOverflow & 53 & 12 & 9 & 32 \\
    All other sources & 66 & 27 & 21 & 18 \\
    \midrule
    \textbf{Total} & \textbf{2615} & \textbf{1029} & \textbf{836} & \textbf{750} \\
    \bottomrule
    \end{tabular}
\end{table}


Table~\ref{tab:ams-breakdown} shows the distribution of statements across AMS subject classifications.
Other notable areas beyond number theory and combinatorics include quantum theory, linear algebra, and algebraic geometry.

\begin{table}[h]
    \caption{Top 10 AMS subject classifications by number of statements in Formal Conjectures.}
    \label{tab:ams-breakdown}
    \centering
    \begin{tabular}{lr}
    \toprule
    \textbf{AMS Subject Classification} & \textbf{Statements} \\
    \midrule
    Number theory & 1545 \\
    Combinatorics & 923 \\
    Quantum theory & 174 \\
    Linear and multilinear algebra; matrix theory & 140 \\
    Information and communication, circuits & 96 \\
    Convex and discrete geometry & 93 \\
    Algebraic geometry & 80 \\
    Geometry & 56 \\
    Field theory and polynomials & 53 \\
    Group theory and generalizations & 42 \\
    \bottomrule
    \end{tabular}
\end{table}


\subsection{The \texorpdfstring{\lean{answer(sorry)}}{answer(sorry)} Mechanism in Detail}\label{sec:app_answer_sorry}

Problems requiring a solution are formalized using \lean{answer( )}, a custom Lean term elaborator.
Concretely, a problem like \enquote{What is the smallest $n$ such that $P(n)$?} is formalized as:
\begin{leancode}
theorem smallest_n : IsLeast {n : ℕ | P n} answer(sorry) := by
  sorry
\end{leancode}
As usual, \lean{sorry} is a placeholder which must be filled in order to claim that a problem is solved, but now there are two of them.
If the solution is known, then \lean{answer(sorry)} is replaced in the statement with \lean{answer(42)}.
This mechanism allows our statements to include computational conjectures where discovering the answer is the primary challenge and verification is routine.

Crucially, while providing a proof simply requires finding any way to replace \lean{:= sorry} with a valid term, providing an answer is a fundamentally different task: it requires evaluating mathematical meaning to determine the correct value.
Lean's type checker can verify that a proposed answer leads to a valid proof, but it cannot judge whether the answer itself is mathematically meaningful---that remains a job for human mathematicians or AI systems with genuine mathematical reasoning capabilities.

There are a number of advantages to using the \lean{answer( )} wrapper over directly writing a solution or a bare \lean{sorry} placeholder:
\begin{itemize}
\item It makes clear which part of the formalization corresponds to the solution; this allows the solution to be masked from \lean{research solved} statements in order to evaluate agents. This can either be done by inspecting the metadata inserted in the Lean term, or more pragmatically, by a simple regex replacement.
\item
    It provides some protection against Lean's flexible elaboration; a problem stated as \lean{1 / 2 * 2 = sorry} could be solved by replacing the \lean{sorry} with either \lean{(0 : ℕ)} or \lean{(1 : ℚ)},
    as by default, Lean only infers the type of an equality after analyzing both sides.
    This is clearly undesirable; we do not want the meaning of our questions to change depending on the answer provided!
    Stating the problem instead as \lean{1 / 2 * 2 = answer(sorry)} protects against this, as the \lean{answer} elaborator postpones analysis which forces Lean to decide the type of the equality before seeing the solution.
\item It forces the statement to be written in such a way that does not pre-assume what the solution is.
\end{itemize}

\subsubsection{Propositional Solutions}

A particularly common use case arises when the truth value of a proposition $P$ is itself unknown, i.e. when the problem can be phrased as a question: "Is it true that...".
In this case, \lean{answer(sorry)} serves as a propositional placeholder that should be replaced with either \lean{True} or \lean{False}:
\begin{leancode}
theorem is_P_true : answer(sorry) ↔ P := by
  sorry
\end{leancode}
Replacing \lean{answer(sorry)} with \lean{True} amounts to conjecturing that $P$ holds (and the solver must prove $P$),
while replacing it with \lean{False} amounts to conjecturing that $P$ fails (and the solver must disprove $P$).
This pattern is widely used throughout the repository for open problems where even the expected truth value is uncertain.

For the convenience of evaluating solver systems that are capable of filling proofs only,
when \lean{answer(sorry)} is detected to be used propositionally in this way, it elaborates to \lean{True} under its default settings.
This permits these systems to decide the proposition by either proving or disproving the statement, rather than by filling the placeholder.


\subsection{The \enquote{for Mathlib} Pattern}\label{sec:app_for_mathlib}

It is typical when undertaking formalization projects to discover missing (and often trivial) results that seem suitable for Mathlib.
Contributing these results to Mathlib is of course the natural resolution, but this has a number of downsides:
\begin{itemize}
\item  New conjectures that depend on these results find themselves at the end of a long dependency chain; the missing results must undergo Mathlib review, the merged change must land in a monthly Mathlib release, and the Formal Conjectures repository must be upgraded to that release
\item Some results may be too specialized for Mathlib, and a Mathlib contribution would require generalizing these immediately.
\item When the results entail a whole new mathematical development, it may take time for the Mathlib community to find the right abstractions and evolve said development into the right shape. Bypassing the dependency chain mentioned above allows much of this evolution to happen more rapidly outside of Mathlib.
\end{itemize}

The Lean community has widely adopted the pattern of creating a \texttt{ForMathlib} directory as a staging ground for such results, the oldest example of which being \href{https://github.com/leanprover-community/lean-liquid/tree/master/src/for\_Mathlib}{the \texttt{for\_mathlib} directory} in the \enquote{liquid tensor experiment}\footnote{\url{https://github.com/leanprover-community/lean-liquid}}.
The formal-conjectures project's \texttt{FormalConjecturesForMathlib} directory follows this pattern, and contains 88 files with results such as natural and logarithmic density of sets of natural numbers (with proofs of basic properties like monotonicity and the density of even numbers), arithmetic progressions and AP-free sets, hypergraph Ramsey numbers, Turing machines and busy beaver halting numbers, perfect powers with a decidability instance, VC dimension in abelian groups, Latin squares and transversals, and a library of graph invariants including the Wiener and Szeged indices, the Lov\'{a}sz theta function, and the Havel--Hakimi residue.
Over time, these results are asynchronously contributed upstream to Mathlib.

\subsection{Misformalization Taxonomy}\label{sec:app_misform_taxonomy}

Misformalizations can be categorized roughly as follows, with the level shown in parentheses.\footnote{Pull request (PR) numbers refer to the repository at \url{https://github.com/google-deepmind/formal-conjectures}; Erd\H{o}s Problem discussions are at \url{https://www.erdosproblems.com}.}
\begin{itemize}
    \item \textbf{Syntactic (1)}: where the formalizer misses a subtlety in the way Lean parses a piece of syntax. For example, this could be missing brackets changing the meaning of an expression in unexpected ways; see PR~\#1338.
    \item \textbf{Semantic (1)}: where the formalizer misses a subtlety in the way that Lean represents types or operators. For example, the natural numbers in Lean contain $0$ and have truncated subtraction; see PRs~\#349, \#1259, \#1262.
    \item \textbf{Misrepresentation (1)}: where the formalizer formally misrepresents aspects of the informal statement during translation. For example, the use of incorrect quantifiers, the use of incorrect logical connectives, or the failure to include hypotheses which are present in the source text; see PRs~\#1164, \#1156.
    \item \textbf{Implicit conventions (2)}: where the source text assumes domain expertise and does not explicitly include certain hypotheses or conventions. For example, in analytic number theory, questions about the magnitude of a set of numbers are often implicitly asymptotic; see PRs~\#2136, \#1151.
    \item \textbf{Reporting (3)}: where the statement changes in the literature through repetition, usually as the result of typos; see Erd\H{o}s Problem~918 discussion.
    \item \textbf{Mathematical (3)}: where the source text is clearly not as intended or contains genuine mathematical errors; see Erd\H{o}s Problem~728 discussion.
\end{itemize}

\begin{table}[H]
    \centering
    \caption{Number and proportion of misformalizations across categories.}
    \label{tab:misform_table}
    \begin{tabular}{llc}
        \textbf{Level} & \textbf{Category} & \textbf{Number of misformalizations} \\
        \toprule
        \multirow{3}{*}{Translation} & Syntactic & 9 (3.09\%)\\
        \cmidrule{2-3}
        & Semantic & 103 (35.40\%)\\
        \cmidrule{2-3}
        & Misrepresentation & 140 (48.11\%) \\
        \midrule
        Underspecified & Implicit conventions & 23 (7.90\%) \\
        \midrule
        \multirow{2}{*}{Source} & Reporting & 15 (5.15\%)\\
        \cmidrule{2-3}
        & Mathematical & 1 (0.34\%) \\
        \bottomrule
    \end{tabular}
\end{table}

\subsection{Misformalization Examples}\label{sec:app_misformalization_examples}

This section contains the misformalization examples, mentioned in Section~\ref{sec:misform}. 

\subsubsection{Syntactical Errors}

An example of a syntactical error involving the correction of misplaced parentheses is shown in Figure \ref{fig:misform_syntactic_erdos1054}.

\begin{figure}[h]
    \centering
    \begin{minted}{diff}
  /-- Let $f(n)$ be the minimal integer $m$ such that $n$ is the sum of the 
  $k$ smallest divisors of $m$ for some $k\geq 1$. Is it true that 
  $\limsup f(n)/n=\infty$? -/
  @[category research open, AMS 11]
  theorem erdos_1054.parts.iii : (∃ (A : Set ℕ), A.HasDensity 1 ∧
-      atTop.limsup (fun n ↦ (f n : EReal) / n) = ⊤) ↔ answer(sorry) := by
+      atTop.limsup (fun n ↦ (f n : EReal) / n = ⊤)) ↔ answer(sorry) := by
  sorry
    \end{minted}
    \caption{Code diff showing an example of a syntactical error, and the correction of parentheses.}
    \label{fig:misform_syntactic_erdos1054}
\end{figure}

\subsubsection{Semantic Errors}

Figures 
\ref{fig:misform_semantic_coeffs}, \ref{fig:misform_semantic_growth} and \ref{fig:misform_semantic_covering} illustrate various semantic errors encountered during formalization.

\begin{figure}[h]
    \centering
    \begin{minted}{diff}
  /-- `Polynomial.HasOddCoeffs f` means that all coefficients of 
  `f : Polynomial ℤ` are odd. -/
  def Polynomial.HasOddCoeffs (f : Polynomial ℤ) : Prop :=
-   ∀ i : ℕ, Odd (f.coeff i)
+   ∀ i ∈ f.support, Odd (f.coeff i)
    \end{minted}
    \caption{Code diff showing an example of a semantic error, with the fix requiring quantification over the support of the polynomial.}
    \label{fig:misform_semantic_coeffs}
\end{figure}

\begin{figure}[h]
    \centering
    \begin{minted}{diff}
  /-- A group `HasPolynomialGrowth` if there exists a finite generating set 
    such that the growth function is bounded above by a polynomial. -/
  def HasPolynomialGrowth (G : Type*) [Group G] : Prop :=
    ∃ (S : Set G), Set.Finite S ∧ Subgroup.closure S = ⊤ ∧
      ∃ (C : ℝ) (d : ℕ), C > 0 ∧
-     ∀ n : ℕ, (GrowthFunction S n : ℝ) ≤ C * (n : ℝ) ^ d
+     ∀ n > 0, (GrowthFunction S n : ℝ) ≤ C * (n : ℝ) ^ d
    \end{minted}
    \caption{Code diff showing an example of a semantic error, where the formalization had failed to account for behavior at the natural number $0$.}
    \label{fig:misform_semantic_growth}
\end{figure}

\begin{figure}
    \centering
    \begin{minted}{diff}
+ /-- An exact covering of a group `G` is a finite collection of subgroups 
+ `{H_1, ..., H_k}` and representative `{g_1, ..., g_k}` such that the 
+ cosets `g_iH_i` are pairwise disjoint and their union covers `G`.
+ 
+ Note that this differs from `Partition (α := Subgroup G)` because the 
+ covering condition there invokes `Subgroup.sup` which is subgroup generation 
+ and thus stronger than union. This definition is easier to use in this 
+ context than the alternative `Partition (α := Set G)`, which lacks
+ subgroup definitions such as `Subgroup.index`. -/
+ structure Group.ExactCovering (G : Type*) [Group G] (ι : Type*) [Fintype ι] where  
+   parts : ι → Subgroup G
+   reps : ι → G
+   nonempty (i : ι) : (parts i : Set G).Nonempty
+   disjoint : (Set.univ (α := ι)).PairwiseDisjoint 
+     fun (i : ι) ↦ reps i • (parts i : Set G)
+   covers : ⋃ i, reps i • (parts i : Set G) = Set.univ

  /--
  If `G` is a group then can there exist an exact covering of `G` by more than 
  one cosets of different sizes? (i.e. each element is contained in exactly 
  one of the cosets.)
  -/
  @[category research open, AMS 20]
  theorem erdos_274 (G : Type*) [Group G] (hG : 1 < ENat.card G) :
-     (∃ (P : Partition (⊤ : Subgroup G)),
-       1 < P.parts.ncard ∧
-       (∀ A ∈ P.parts, ∃ᵉ (s : G) (H : Subgroup G), s • (H : Set G) = A) ∧
-       P.parts.Pairwise fun A B ↦ #A ≠ #B) ↔ answer(sorry) := by
+     (∃ (ι : Type*) (_ : Fintype ι) (P : Group.ExactCovering G ι),
+       1 < Fintype.card ι ∧ (Set.range P.parts).Pairwise fun A B ↦ #A ≠ #B) ↔ 
+         answer(sorry) := by
  sorry
    \end{minted}
    \caption{Code diff showing an example of a semantic error, where the formalization missed the subtlety that \lean{Partition} invokes \lean{sup} for the covering condition, which means something different for subgroups than for sets.}
    \label{fig:misform_semantic_covering}
\end{figure}

\clearpage 

\subsubsection{Misrepresentation Errors} 

Figure \ref{fig:misform_misrep_giuga} provides an example of a misrepresentation error where a hypothesis from the informal text was omitted, and Figure \ref{fig:misform_misrep_quantifier} shows an error involving incorrect quantification.

\begin{figure}[h]
    \centering
    \begin{minted}{diff}
  /--
- A (weak) Giuga number is a number $n$ such that
+ A (weak) Giuga number is a composite number $n$ such that
  $$\sum_{i=1}^{n - 1}i^{\varphi(n)} \equiv -1\pmod{n}$$.
  -/
  def IsWeakGiuga (n : ℕ) : Prop :=
-     2 ≤ n ∧ ¬ n.Prime ∧ n ∣ 1 + ∑ i ∈ Finset.Ioo 0 n, i ^ φ n
+     n.Composite ∧ n ∣ 1 + ∑ i ∈ Finset.Ioo 0 n, i ^ φ n
    \end{minted}
    \caption{Code diff showing an example of a misrepresentation error, where the formalization was missing a hypothesis that was present in the informal text. }
    \label{fig:misform_misrep_giuga}
\end{figure}

\begin{figure}[h]
    \centering
    \begin{minted}{diff}
  @[category research open, AMS 11]
  theorem erdos_944 :
-     (∀ k ≥ 4, ∀ r ≥ 1, ∃ (G : SimpleGraph V), G.IsErdos944 k r) ↔ 
-       answer(sorry) := by
+     (∀ k ≥ 4, ∀ r ≥ 1, ∃ (V : Type u) (G : SimpleGraph V), G.IsErdos944 k r) ↔ 
+       answer(sorry) := by
  sorry
    \end{minted}
    \caption{Code diff showing an example of a misrepresentation error, where the vertex type \lean{V} was incorrectly quantified. In the old version \lean{V} was a section variable and thus appeared out of the scope of the existential quantifier.}
    \label{fig:misform_misrep_quantifier}
\end{figure}

\subsubsection{Implicit Conventions}

Examples of errors stemming from implicit conventions in the source material can be seen in Figures \ref{fig:misform_implicit_asymptotic} and \ref{fig:misform_implicit_induced}.

\begin{figure}[h]
    \centering
    \begin{minted}{diff}
  @[category research open, AMS 11]
  theorem erdos_510 :
      answer(sorry) ↔ ∃ (c : ℝ) (hc : 0 < c),
-       ∀ N > 0, ∀ (A : Finset ℕ), 0 ∉ A → #A = N →
+       ∀ᶠ N in atTop, ∀ (A : Finset ℕ), 0 ∉ A → #A = N →
        ∃ θ, ∑ n ∈ A, cos (n * θ) < -c * sqrt N := by
    sorry
    \end{minted}
    \caption{Code diff showing an example of an implicit convention, where the source\protect\footnotemark\ implicitly assumes that the inequality holds for sufficiently large $N$. This is a common implicit convention across analytic number theory where one is typically interested in the asymptotic behavior of functions or sequences of natural numbers.}
    \label{fig:misform_implicit_asymptotic}
\end{figure}
\footnotetext{\url{https://www.erdosproblems.com/510}}
\newpage
\begin{figure}[h]
    \centering
    \begin{minted}{diff}
  @[category research open, AMS 5]
  theorem erdos_128 :
-     ((∀ (G' : G.Subgraph) [Fintype G'.verts] [Fintype G'.edgeSet],
-         letI n := Fintype.card V;
-         2 * G'.verts.toFinset.card ≥ n →
-         50 * G'.edgeSet.toFinset.card > n^2) → ¬ (G.CliqueFree 3))
+     ((∀ (V' : Set V),
+       2 * V'.ncard ≥ Fintype.card V →
+         50 * (G.induce V').edgeSet.ncard > Fintype.card V ^ 2) → ¬(G.CliqueFree 3))
    ↔ answer(sorry) := by
  sorry
    \end{minted}
    \caption{Code diff showing an example of an implicit convention, where the source\protect\footnotemark\ wrote ``subgraph" rather than ``induced subgraph". The comments\protect\footnotemark\ demonstrate that the problem is trivial if one does not restrict to induced subgraphs, which is why we class this as implicit convention. It might also be considered to be a reporting error since the original paper explicitly writes ``induced subgraph", \cite[p. 344]{Erdos93}.}
    \label{fig:misform_implicit_induced}
\end{figure}
\addtocounter{footnote}{-1}\footnotetext{\url{https://www.erdosproblems.com/history/128}}
\addtocounter{footnote}{1}\footnotetext{\url{https://www.erdosproblems.com/forum/thread/128}}
 
\subsubsection{Reporting Errors}

Figure \ref{fig:misform_reporting_erdos918} demonstrates a reporting error caused by ambiguities in the historical source material regarding induced subgraphs.

\begin{figure}[H]
    \centering
    \begin{minted}{diff}
  /-- Is there a graph with $\aleph_2$ vertices and chromatic number $\aleph_2$ such that every
  subgraph on $\aleph_1$ vertices has chromatic number $\leq\aleph_0$? -/
  -- Formalisation note: source material [ErHa68b] uses only induced subgraphs
  @[category research open, AMS 5]
  theorem erdos_918.parts.i :
      (∃ (V : Type u) (G : SimpleGraph V), #V = ℵ_ 2 ∧ G.chromaticCardinal = ℵ_ 2 ∧
-       ∀ (W : Set V) (_ : #W = ℵ₁), (G.induce W).chromaticCardinal = ℵ₀) ↔
+       ∀ (W : Set V) (_ : #W = ℵ₁), (G.induce W).chromaticCardinal ≤ ℵ₀) ↔
      answer(sorry) := by
    sorry
    \end{minted}
    \caption{Code diff showing an example of a reporting error between an internal draft formalization of the problem. This problem appears in \cite{EH68} and \cite{Erdos69}. In the former citation the explicit use of $\leq \aleph_0$ is used, while this does not appear in the latter citation. This ambiguity was only clarified post draft formalization with the discovery of trivial solutions -- in the (non-induced) subgraph case by the website comments\protect\footnotemark\ and by AlphaProof in the induced subgraph case. Note also that \cite{EH68} explicitly writes \emph{induced} subgraph while \cite{Erdos69} and the website do not.}
    \label{fig:misform_reporting_erdos918}
\end{figure}
\footnotetext{\url{https://www.erdosproblems.com/forum/thread/918}}
